% BHU PYQ -- Transcribed & Corrected
% Paper: BOTMJ41: Gymnosperms and Paleobotany
% Exam: B.Sc. (Hons.) Semester IV (NEP) Examination, 2025-26
% Ref Code: D26(232)-300
% Category: Major / Minor
% Department: Botany

\documentclass[11pt,a4paper]{article}
\usepackage[a4paper,top=2.5cm,bottom=2.5cm,left=2.8cm,right=2.8cm,headheight=14pt]{geometry}
\usepackage{amsmath,amssymb}
\usepackage[T1]{fontenc}
\usepackage[utf8]{inputenc}
\usepackage{lmodern,microtype,fancyhdr,enumitem,hyperref,lastpage,parskip}

\hypersetup{
    colorlinks=true,
    urlcolor=black,
    linkcolor=black,
    pdftitle={BOTMJ41: Gymnosperms and Paleobotany -- BHU Sem IV 2025-26 (NEP)}
}

\pagestyle{fancy}
\fancyhf{}
\renewcommand{\headrulewidth}{0.4pt}
\renewcommand{\footrulewidth}{0.4pt}
\fancyhead[L]{\small \textit{Banaras Hindu University}}
\fancyhead[R]{\small \textit{BOTMJ41: Gymnosperms \& Paleobotany (2025-26)}}
\fancyfoot[C]{\small Page \thepage\ of \pageref{LastPage}}

\newcommand{\pts}[1]{\hfill \textit{[#1]}}

\begin{document}

\begin{center}
  {\large\bfseries BANARAS HINDU UNIVERSITY}\\[2pt]
  {\normalsize B.Sc. (Hons) Semester IV Examination, 2025-26 (NEP)}\\[6pt]
  \rule{0.8\linewidth}{0.5pt}\\[6pt]
  {\large\bfseries BOTANY}\\[4pt]
  {\large\bfseries Paper Code: BOTMJ41 : Gymnosperms and Paleobotany}\\[2pt]
  {\small Ref. No.: D26(232)-300}\\[4pt]
  {\normalsize Time: 2$\frac{1}{2}$ Hours\hfill Maximum Marks: 70}
\end{center}

\rule{\linewidth}{0.4pt}

\smallskip

\noindent \textit{Instructions: Answer five questions including Question 1 which is compulsory. The figures in the right-hand margin indicate marks.}

\bigskip

\noindent \textbf{Question 1.} Answer the following questions in not more than 100 words each: \pts{7 $\times$ 2 = 14}
\begin{enumerate}[label=(\alph*)]
  \item Why are the leaves of \textit{Cycas} described as pinnately compound but fern-like in appearance?
  \item Describe the salient features of Progymnosperm.
  \item Why are the leaves of \textit{Ephedra} reduced and what is the role of its green, jointed stems?
  \item Give an account of the affinities of Pteridospermales with Cycadophyta and ferns.
  \item Why are the seeds of \textit{Pinus} winged and how are these wings derived?
  \item Describe the anatomical features of \textit{Lyginopteris}.
  \item What is the significance of the presence of resin canals for survival in \textit{Pinus}?
\end{enumerate}

\bigskip

\noindent \textbf{Question 2.} Describe the geological time scale, outlining its major divisions, key events in each era and its significance in understanding the evolution of plant life. \pts{14}

\bigskip

\noindent \textbf{Question 3.} Write short notes on \textit{any two} of the following: \pts{2 $\times$ 7 = 14}
\begin{enumerate}[label=(\alph*)]
  \item Reproductive parts of \textit{Lyginopteris}
  \item Pteridospermales
  \item Classification of gymnosperms.
\end{enumerate}

\bigskip

\noindent \textbf{Question 4.} Give a comprehensive account of the types of fossils and their importance, with suitable diagrams. \pts{14}

\bigskip

\noindent \textbf{Question 5.} Write short notes on \textit{any two} of the following: \pts{2 $\times$ 7 = 14}
\begin{enumerate}[label=(\alph*)]
  \item Reproduction in \textit{Pinus}
  \item Difference between young and old stems of \textit{Ephedra}
  \item Structure and adaptations of \textit{Pinus} leaf.
\end{enumerate}

\bigskip

\noindent \textbf{Question 6.} Describe the anatomy and reproduction of \textit{Cycas}, highlighting its distinctive features and evolutionary significance. \pts{14}

\bigskip

\noindent \textbf{Question 7.} Write short notes on \textit{any two} of the following: \pts{2 $\times$ 7 = 14}
\begin{enumerate}[label=(\alph*)]
  \item Major geological events in Earth's history
  \item \textit{Glossopteris} leaf
  \item Poikilohydry and homoiohydry.
\end{enumerate}

\bigskip
\begin{center}
  \textbf{***}
\end{center}

\end{document}
